%============================================================
% CHAPITRE 2 — CONTEXTE GÉNÉRAL DU PROJET
%============================================================
\chapter{Contexte général du projet}

\section{Introduction}

MediCab Pro est une solution numérique de type SaaS pensée pour répondre aux
besoins quotidiens des professionnels de santé exerçant en libéral au Maroc.
Sa cible principale regroupe les médecins généralistes et spécialistes, leurs
assistants administratifs ainsi que les gestionnaires de la plateforme. Le périmètre
fonctionnel de la solution couvre l'intégralité du parcours de soin, de la prise
de rendez-vous initiale jusqu'à l'émission des factures, en englobant la tenue du
dossier patient électronique et la rédaction assistée des prescriptions médicales
par intelligence artificielle.

La pertinence de ce projet tient à sa capacité à transformer structurellement
le fonctionnement des cabinets médicaux marocains, en leur proposant une solution
numérique complète, sécurisée et adaptée aux contraintes locales : qualité variable
de la connexion internet, usage quotidien du français et de l'arabe, et respect
des exigences légales de la loi 09-08 relative à la protection des données à
caractère personnel.

\section{Problématique}

L'observation des pratiques actuelles dans les cabinets médicaux privés marocains
révèle que la gestion manuelle reste la norme dans une proportion importante des
structures. Cette réalité engendre plusieurs catégories de dysfonctionnements,
identifiées lors de l'analyse des besoins :

\begin{itemize}
  \item Accès difficile aux archives médicales et risque non négligeable de perte
    ou d'altération des dossiers papier.
  \item Gestion approximative des agendas (rendez-vous qui se chevauchent,
    absences non signalées, absence de rappels automatiques aux patients).
  \item Suivi financier défaillant : paiements non tracés, remboursements
    CNSS/CNOPS non consolidés, impayés mal gérés.
  \item Absence d'intégration des technologies actuelles : Cloud, architectures
    distribuées, intelligence artificielle.
  \item Déploiements applicatifs artisanaux, difficiles à maintenir et exposés
    aux incidents de production.
  \item Manque d'outils adaptés au bilinguisme du marché marocain (français/arabe).
\end{itemize}

\section{Cahier des charges}

\subsection{Introduction}

L'ambition principale de ce projet est de livrer une plateforme numérique complète
permettant aux équipes médicales de prendre en charge l'ensemble des activités
administratives et cliniques d'un cabinet, depuis n'importe quel navigateur web,
sans nécessiter d'installation locale ni d'infrastructure propre.

\subsection{Besoins fonctionnels}

Les fonctionnalités attendues de la plateforme sont organisées en six grands modules :

\subsubsection*{Module 1 — Gestion des Rendez-vous}
\begin{itemize}
  \item Calendrier interactif actualisé en temps réel (vues par jour, semaine et mois).
  \item Création, modification et suppression de rendez-vous avec gestion
    des états (confirmé, urgent, en attente, annulé).
  \item Liste d'attente dynamique et traitement des absences patients.
  \item Système de rappels automatiques par notification.
\end{itemize}

\subsubsection*{Module 2 — Dossier Patient Électronique (DPE)}
\begin{itemize}
  \item Constitution et administration des dossiers patients (identité, antécédents
    médicaux et chirurgicaux, allergies déclarées, traitements en cours).
  \item Suivi des paramètres biologiques : groupe sanguin, poids, taille, avec
    calcul automatique de l'indice de masse corporelle (IMC).
  \item Historique exhaustif des consultations avec diagnostic et montant associé.
  \item Archivage et consultation des documents médicaux avec export au format PDF.
  \item Recherche multicritère par nom, numéro de CIN ou numéro de téléphone.
\end{itemize}

\subsubsection*{Module 3 — Consultations et Ordonnances}
\begin{itemize}
  \item Saisie des consultations avec motif de visite, type, diagnostic retenu et montant.
  \item Élaboration et gestion des ordonnances électroniques avec liste de médicaments.
  \item Impression et export au format PDF des ordonnances et comptes rendus médicaux.
\end{itemize}

\subsubsection*{Module 4 — Facturation et Gestion Financière}
\begin{itemize}
  \item Émission automatique des factures liées aux actes de consultation.
  \item Suivi du statut des paiements (acquittée, en attente, non réglée).
  \item Tableau de bord financier avec indicateurs graphiques interactifs.
  \item Export des données en PDF et en format Excel.
\end{itemize}

\subsubsection*{Module 5 — Assistant IA Médical (MedAgent)}
\begin{itemize}
  \item Interface intelligente permettant la dictée ou la saisie textuelle des
    observations cliniques.
  \item Extraction automatique des informations médicales pertinentes (symptômes,
    diagnostic probable).
  \item Rédaction automatique de comptes rendus médicaux structurés.
  \item Génération automatique d'ordonnances fondée sur le diagnostic établi.
\end{itemize}

\subsubsection*{Module 6 — Administration Multi-Tenant}
\begin{itemize}
  \item Tableau de bord dédié au super-administrateur avec vision globale de
    la plateforme.
  \item Gestion du cycle de vie des cabinets médicaux (création, activation,
    suspension).
  \item Supervision des abonnements et des grilles tarifaires.
  \item Monitoring de la plateforme avec indicateurs de croissance.
\end{itemize}

\subsection{Besoins non fonctionnels}

\begin{table}[H]
\centering
\caption{Exigences non fonctionnelles de MediCab Pro}
\label{tab:nonfonctionnels}
\renewcommand{\arraystretch}{1.4}
\begin{tabularx}{\textwidth}{|>{\bfseries}p{3.5cm}|X|}
\hline
\textbf{Catégorie} & \textbf{Détails} \\
\hline
Disponibilité &
  Taux de disponibilité de 99.9\% assuré par l'orchestration Kubernetes
  et l'auto-scaling horizontal. \\
\hline
Sécurité &
  Protocole HTTPS obligatoire, authentification par JWT, contrôle d'accès
  basé sur les rôles (RBAC), conformité à la loi 09-08 marocaine. \\
\hline
Performance &
  Temps de réponse des API REST inférieur à 300 ms pour 95\% des requêtes.
  Latence P99 sous le seuil de 500 ms. \\
\hline
Scalabilité &
  Architecture multi-tenant, montée en charge horizontale via le HPA
  Kubernetes. Capacité à supporter plus de 100 cabinets simultanés. \\
\hline
Multi-langue &
  Interface disponible en français et en arabe. \\
\hline
Responsivité &
  Interface adaptative compatible avec les postes de travail, tablettes
  et terminaux mobiles. \\
\hline
Maintenabilité &
  Architecture microservices, pipeline CI/CD intégré, analyse continue
  par SonarQube. Couverture de tests supérieure à 80\%. \\
\hline
Mode sombre &
  Prise en charge native du thème sombre (Dark Mode). \\
\hline
\end{tabularx}
\end{table}

\section{Spécification des besoins}

\subsection{Identification des acteurs}

\begin{table}[H]
\centering
\caption{Acteurs du système MediCab Pro et leurs attributions}
\label{tab:acteurs}
\renewcommand{\arraystretch}{1.4}
\begin{tabularx}{\textwidth}{|>{\bfseries}p{3cm}|X|}
\hline
\textbf{Acteur} & \textbf{Rôle et attributions} \\
\hline
Médecin &
  \begin{itemize}[leftmargin=*,noitemsep,topsep=2pt]
    \item Administration de son planning professionnel
    \item Consultation des dossiers patients et de l'historique médical
    \item Création et gestion des consultations
    \item Établissement et impression des prescriptions médicales
    \item Recours à l'assistant intelligent MedAgent
    \item Analyse des indicateurs financiers du cabinet
  \end{itemize} \\
\hline
Secrétaire Médicale &
  \begin{itemize}[leftmargin=*,noitemsep,topsep=2pt]
    \item Organisation des rendez-vous et du planning médical
    \item Accueil et orientation des patients (file d'attente)
    \item Enregistrement des nouveaux patients dans le système
    \item Gestion de la facturation et des encaissements
  \end{itemize} \\
\hline
Super Administrateur &
  \begin{itemize}[leftmargin=*,noitemsep,topsep=2pt]
    \item Supervision générale de la plateforme
    \item Administration des cabinets médicaux (approche multi-tenant)
    \item Gestion des utilisateurs et des abonnements
    \item Supervision des performances et gestion des alertes système
  \end{itemize} \\
\hline
Système IA (MedAgent) &
  \begin{itemize}[leftmargin=*,noitemsep,topsep=2pt]
    \item Traitement et analyse du discours médical
    \item Extraction des entités cliniques pertinentes
    \item Production automatique de comptes rendus et d'ordonnances
  \end{itemize} \\
\hline
\end{tabularx}
\end{table}

\subsection{Conclusion}

Ce chapitre a établi le cadre fonctionnel et technique de MediCab Pro, en
précisant les attentes des différentes catégories d'utilisateurs ainsi que les
contraintes du projet. Le chapitre suivant sera consacré à la phase d'analyse
et de conception, avec la modélisation UML et la présentation détaillée de
l'architecture technique retenue.
